\thispagestyle{empty}
\vspace*{\fill}
\begin{small}\sffamily
\textbf{Manual del usuario de REM}\\
Autor: Ramon Guillamon \textless learntouseai@gmail.com\textgreater\\
Edición 1.0, octubre de 2026.

\medskip\textbf{Licencias.}
\begin{itemize}
\item \textbf{Este libro} se publica bajo la licencia \textbf{Creative Commons Atribución-CompartirIgual 4.0 Internacional} (CC BY-SA 4.0). Puedes copiarlo, redistribuirlo y adaptarlo, incluso con fines comerciales, siempre que reconozcas la autoría, indiques si has hecho cambios y distribuyas tus obras derivadas con la misma licencia. Texto completo: \url{https://creativecommons.org/licenses/by-sa/4.0/deed.es}.
\item \textbf{El código fuente de REM} que acompaña al libro (el repositorio) se publica bajo la licencia \textbf{MIT}: puedes usarlo, modificarlo y redistribuirlo, también con fines comerciales, conservando el aviso de copyright.
\end{itemize}

\medskip\textbf{Datos ficticios.} Todas las capturas de pantalla de este libro se han hecho con datos inventados. Los nombres de personas y empresas, los logotipos, los NIF (de ceros), los teléfonos, los importes, los correos (en el dominio reservado \texttt{.invalid}) y los IBAN no corresponden a ninguna persona ni entidad reales. Cualquier parecido con la realidad es casual.

\medskip\textbf{Aviso.} REM no es software de facturación certificado (VeriFactu) ni ofrece asesoramiento fiscal, legal o contable. Las cifras y los modelos fiscales que muestra son una ayuda; la responsabilidad de lo que presentes es tuya y de tu asesor. El libro describe un prototipo en evolución: el canal con clientes (REM Clientes) no ha pasado una auditoría de seguridad independiente.

\medskip\textbf{Marcas.} Apple, macOS, Telegram, WhatsApp, Claude y otros nombres citados pertenecen a sus respectivos titulares y se mencionan solo a efectos descriptivos.
\end{small}
\vspace*{\fill}
\cleardoublepage
